\documentclass{article}
\usepackage[oldstyle]{ebgaramond}
\usepackage{propositions}

%% The digit widening with proportional (oldstyle) figures.  Computer Modern cannot
%% exercise this: all of its digits have the same width, so any choice of
%% substitute digit gives the same answer and a wrong choice cannot show up.
%% EB Garamond's oldstyle figures are genuinely proportional --- 0 is the
%% widest and 1 is 44% narrower --- so here the substitute has to be the widest
%% digit and nothing else will do.  (Requires ebgaramond.)

%% \propwidest is package-internal (\__props_widest:n).  A test may reach for
%% it anyway; this shim keeps the prose below readable without pretending the
%% command is public.
\ExplSyntaxOn
\NewDocumentCommand \widen { m } { \__props_widest:n {#1} }
\ExplSyntaxOff

\newlength\wa \newlength\wb
\newcommand\cmp[1]{%
  \settowidth\wa{\rmfamily\normalsize(#1)}%
  \settowidth\wb{\rmfamily\normalsize\widen{(#1)}}%
  \texttt{(#1)}: actual \the\wa, reserved \the\wb,
  \ifdim\wa>\wb \textbf{OVERFLOW by \the\dimexpr\wa-\wb\relax}%
  \else slack \the\dimexpr\wb-\wa\relax\fi\par}

\begin{document}

\section{Nothing may overflow}

The reserved width is what \verb|fitmargin| sets aside for the label.  Every
line below must read \emph{slack}: a label wider than the space reserved for it
overflows its box and pushes the first line of the item body sideways, which is
the very jitter the widening exists to remove.  Reserving for \texttt{(88)}
rather than for the widest variant fails here on every label containing a
\texttt{0}, \texttt{4}, \texttt{6}, \texttt{9} or \texttt{7}.

{\ttfamily\footnotesize
\cmp{1}\cmp{5}\cmp{8}\cmp{9}\cmp{0}
\cmp{11}\cmp{99}\cmp{00}\cmp{40}\cmp{88}
\cmp{100}\cmp{999}\cmp{000}
}

Exactly one line in each group must show zero slack --- the all-widest-digit
label itself, which is what was measured.

\section{Equal digit counts, equal margins}

\noindent\rule{\textwidth}{0.3pt}
\setcounter{numpropi}{10}
\begin{prop}[style=fitmargin]
  \item Label (11): the two narrowest digits, and the loosest gap.
\end{prop}
\setcounter{numpropi}{98}
\begin{prop}[style=fitmargin]
  \item Label (99): must begin at exactly the same place as the item above,
  even though its label is 3.5pt wider.
\end{prop}
\setcounter{numpropi}{999}
\begin{prop}[style=fitmargin]
  \item Label (1000): four digits, so this one is indented further.
\end{prop}
\noindent\rule{\textwidth}{0.3pt}
\setcounter{numpropi}{0}

The bodies of the first two items must start at the same horizontal position;
the third starts further right.  Reserving for \texttt{(88)} instead leaves the
second one 0.28pt out.

\section{Non-digits are untouched}

\verb|\__props_widest:n| of \textbf{(3.147)} is \widen{\textbf{(3.147)}}, still
bold and with the point still in place; of \textsc{Prop}~1a it is
\widen{\textsc{Prop}~1a}, still small caps with the space intact; and of a
label with no digits at all, \emph{Nihilism}, it is \widen{\emph{Nihilism}}
--- unchanged, and measured no times at all.

\end{document}
